(1) 圆半径 \(R\)，周长 \(L=2\pi R\)，接收点切向速度 \(v=\Omega R\)。从实验室系看，顺向光需追及移动点，故 \(t_+=L/(c-v)\)，逆向与点相向，\(t_-=L/(c+v)\)（小旋转忽略高阶路径弯曲）。差值为
\[
\Delta t=t_+-t_-=\frac{2Lv}{c^2-v^2}
=\frac{4\Omega A}{c^2}+O(\Omega^3),\quad A=\pi R^2.
\]
在小 \(v/c\) 限制下得到一阶面积律，若忽略探测点的旋转而两次都取 \(L/c\)，便错误抹去Sagnac信号。

(2) 对任意平面闭合环的一阶传播时差，速度场 \(\boldsymbol v=\boldsymbol\Omega\times\boldsymbol r\) 沿线路的对流修正使
\[
\Delta t=\frac2{c^2}\oint\boldsymbol v\cdot d\boldsymbol\ell
=\frac2{c^2}\oint(\boldsymbol\Omega\times\boldsymbol r)\cdot d\boldsymbol\ell.
\]
Stokes定理及 \(\nabla\times(\boldsymbol\Omega\times\boldsymbol r)=2\boldsymbol\Omega\) 给
\[
\boxed{\Delta t=4\boldsymbol\Omega\cdot\boldsymbol A/c^2}.
\]
面积是有向面积矢量，因此翻转环法向或旋转方向会改变信号符号；同面积不同形状的闭合路径首阶时差相同。

(3) 两束回到干涉端口的角频率为 \(\omega_0=2\pi c/\lambda\)，相位差
\[
\boxed{\Delta\phi=\omega_0\Delta t
=\frac{8\pi\,\boldsymbol\Omega\cdot\boldsymbol A}{\lambda c}}.
\]
同一方向连续绕行 \(m\) 圈，累积有向面积为 \(m\boldsymbol A\)，故信号乘 \(m\)。此处 \(\lambda\) 为真空波长；若以角频率与普通频率不分，或漏掉 \(2\pi\)，会产生数值错误。

(4) 单圈最优干涉工作点的相位标准差散粒极限为 \(\delta\phi=1/\sqrt N\)，由灵敏度导数 \(|\partial\phi/\partial\Omega|=8\pi A/(\lambda c)\) 得
\[
\boxed{\delta\Omega_{\rm shot}=\frac{\lambda c}{8\pi A\sqrt N}}.
\]
多圈时若入射及最终探测光子数仍固定为 \(N\)，公式分母再乘 \(m\)；现实中额外损耗会降低 \(N\)，镜面漂移和震动造成非散粒噪声，且高速旋转需用完整相对论传播而非此首阶展开。
